\section{The \crystal Model}
\label{sec:model}

\subsection{Scope and Rated Attribute}
\label{sec:model-scope}
\crystal models the \emph{network-attributable} component of \qoe for
\vpcc streaming over reliable \quic: the degradation in perceived
temporal smoothness caused by delay and loss, relative to an
impairment-free presentation of the same content. It does not model
compression fidelity, which is fixed per sequence and reported through
\pcqm as a stimulus descriptor, and it does not model overall \mos as a
combination of fidelity and smoothness. The rated attribute in the
subjective study is smoothness of playback. One quantity in the model
is fitted to ratings ($\alpha$), one is measured on the clean slice
($\qc$), and everything else is a transport fact, a client property or
a session counter (Table~\ref{tab:notation}).

\begin{table}[t]
  \caption{Notation. Values in parentheses are those of the campaign
  client of Section~\ref{sec:system}.}
  \label{tab:notation}
  \centering
  \footnotesize
  \setlength{\tabcolsep}{4pt}
  \begin{tabular}{@{}lll@{}}
    \toprule
    Symbol & Meaning & Source \\
    \midrule
    $G$, $F$ & \gof size (frames), frame rate & encoder \\
    $T_G = G/F$ & content-clock period of one delivery unit & derived (\SI{133.3}{\milli\second} at $G=4$) \\
    $f$ & client render-frame length & client, logged per session (\SI{30}{\milli\second}) \\
    $B_{\mathrm{eff}} = f(\lfloor T_G/f \rfloor - 1)$ & per-unit budget after frame quantization & derived (\SI{90}{\milli\second}) \\
    $d$ & added one-way delay & network condition \\
    $\tau_{\mathrm{dl}}$ & clean transfer time of a unit & pipeline log at $P=0$ ($\approx$35--\SI{39}{\milli\second}) \\
    $S(d) = B_{\mathrm{eff}} - \tau_{\mathrm{dl}} - d$ & slack: margin left for loss recovery & derived \\
    $P$ & packet loss rate & network condition \\
    $N_p$ & packets per unit & log ($\approx$471--517) \\
    $\mu = P N_p$ & expected losses per unit & derived \\
    $r(d) = r_0 + \kappa d$ & effective cost of one loss, incl.\ window cut & server transport logs \\
    $K(d) = \lfloor S(d)/r(d) \rfloor$ & losses a unit can absorb & derived \\
    $H$ & decoder lead over the content clock & client ($\approx$\SI{550}{\milli\second}, $\approx$4 units) \\
    $\tlf$ & skip fraction: (skipped + jumped) / due & session counters \\
    $\phi$ & freeze fraction: units shown late / due & session counters \\
    $\qc$ & clean-reference ceiling, $(\mosm_{A,c} - 1)/4$ & measured, Slice~A \\
    $\alpha$ & perceptual sensitivity to lost playout & fitted to ratings \\
    \bottomrule
  \end{tabular}
\end{table}

\subsection{Delivery Unit and Loss Recovery}
\label{sec:model-unit}
Two properties of the stream and the transport determine how a network
impairment reaches the viewer. The first is the delivery unit. \vpcc,
the principal component of ISO/IEC 23090-5~\cite{v3c,vpccgraziosi},
projects a dynamic point cloud into geometry, attribute and occupancy
atlases coded as ordinary video with \gls{hevc}~\cite{hevc,tmc2}, and so
inherits the video codec's prediction structure: frames are coded in
\glspl{gof}, either \gls{ra} with I-, P- and B-frames and a \gof of 32,
or \gls{ldp} with I- and P-frames and \gofs of 16 or
4~\cite{vpccctc,hmctc}. A \gof must be received in full before any of
its frames can be reconstructed, and in a streaming deployment it
carries a presentation deadline. The \gof is therefore the
\emph{delivery unit}, and its size sets the budget within which a
disturbance must be absorbed.

The second is loss recovery. Each \gof is fetched as one HTTP object
over \gls{http3}~\cite{rfc9114} on one \quic stream~\cite{rfc9000}.
\quic declares a packet lost either when a packet at least
$k_{\mathrm{PacketThreshold}}=3$ later is acknowledged or, if no later
acknowledgment arrives, when a \gls{pto} expires~\cite{rfc9002}:
\begin{equation}
  \mathrm{PTO} = \mathrm{sRTT} + \max(4\,\mathrm{RTT}_{\mathrm{var}},
  k_{\mathrm{granularity}}) + \mathrm{max\_ack\_delay}.
  \label{eq:pto}
\end{equation}
A loss therefore costs at least one round trip, and each recovery is
followed by a congestion-window reduction that persists into later
transfers. Because \quic is reliable, the decoder never receives an
incomplete \gof: loss becomes \emph{delay}, and delay becomes a
\emph{freeze} or a \emph{skip} when the unit misses its presentation
time. Our measurements show the two impairments leave distinct
signatures, added latency shifting the delivery-time distribution and
loss expanding its tail~\cite{vquest2026}, and that the per-\gof
deadline-miss rate correlates only partially with subjective quality
($\rho = -0.34$) because a missed \gof of 4 frames and one of 16
frames remove different durations of content~\cite{mosaic2026}.

\subsection{Layer 1: One Perceptual Channel}
\label{sec:model-layer1}

\subsubsection{Premises and claim}
Layer~1 rests on two premises. Under \emph{reliable transport}, a lost
packet never corrupts a decoded unit; it delays the unit's arrival.
Under \emph{deadline-enforced buffered playout}
(Section~\ref{sec:sys-edge}), the client keeps one request in flight
and feeds a decoder that holds about $H$ of content ahead of the
content clock; each unit has a fixed presentation time, a unit arriving
slightly late is shown after a brief freeze, and a unit arriving more
than one period late is skipped. Delay is invisible while the per-unit
cycle keeps pace with the content clock. From the session counters of
the client we define
\begin{equation}
  \tlf = \frac{\text{units skipped} + \text{units jumped over}}{\text{units due}},
  \qquad
  \phi = \frac{\text{units shown late}}{\text{units due}}.
  \label{eq:tlf}
\end{equation}

\emph{Claim 1.} Perceived smoothness depends on the network only
through presentation timing, $(\tlf, \phi)$. Delay and loss have no
perceptual channel of their own: at equal $(\tlf, \phi)$, ratings do
not depend on whether the misses came from loss, delay or both, and
impaired sessions with $\tlf \approx 0$ rate like the clean slice.
Claim~1 is the hypothesis the subjective study tests; the rest of the
model takes it as given and explains how $\tlf$ arises.

\subsubsection{When is a unit slow?}
\label{sec:model-slow}
Consider one unit. Its fetch takes
\begin{equation}
  t_f = d + \tau_{\mathrm{dl}} + \sum_{i=1}^{n} r_i(d),
  \label{eq:fetch}
\end{equation}
the added delay, the clean transfer time and the cost of each of its
$n$ loss events. The client notices completion at the next render-frame
boundary and hands the unit on one frame later, so the cycle costs
$c = f\lceil t_f/f \rceil + f$. The client keeps pace while $c \le T_G$,
which, since $\lceil t_f/f \rceil$ is an integer, is equivalent to
\begin{equation}
  t_f \le B_{\mathrm{eff}} = f\left(\left\lfloor \frac{T_G}{f}\right\rfloor - 1\right).
  \label{eq:budget}
\end{equation}
With $T_G = \SI{133.3}{\milli\second}$ and $f = \SI{30}{\milli\second}$,
$B_{\mathrm{eff}} = \SI{90}{\milli\second}$; a continuously polling
client has $B_{\mathrm{eff}} \to T_G$. Subtracting the parts spent
before any loss occurs leaves the slack
\begin{equation}
  S(d) = B_{\mathrm{eff}} - \tau_{\mathrm{dl}} - d.
  \label{eq:slack}
\end{equation}
Delay therefore acts on its own, before any loss: a unit is slow with
zero loss once $d > B_{\mathrm{eff}} - \tau_{\mathrm{dl}}$. With
$\tau_{\mathrm{dl}} \approx$ 35--\SI{39}{\milli\second} this predicts a
pure-delay threshold of 51--\SI{55}{\milli\second}, with no fitted
parameter; the calibration scan of Section~\ref{sec:calibration}
measured 50--\SI{55}{\milli\second}.

The number of loss events per unit is $n \sim \mathrm{Poisson}(\mu)$,
$\mu = P N_p$, the Poisson approximation to $\mathrm{Binomial}(N_p, P)$
at $P \le 1\%$. Each loss costs detection
and retransmission, both scaling with the round-trip time
(Eq.~\eqref{eq:pto}), and a congestion-window cut (NewReno,
$\beta = 0.7$) that an application-limited sender does not regrow
before the next unit, so one loss also slows the units that follow. We
fold all of this into one effective per-loss cost,
\begin{equation}
  r(d) = r_0 + \kappa d,
  \label{eq:recovery}
\end{equation}
with $r_0$ and $\kappa$ estimated from the server transport logs, not
from ratings. Because the window effect spans units, $r$ is an
effective quantity and this stage is a per-unit approximation of a
process with short memory. With $K(d) = \lfloor S(d)/r(d) \rfloor$
losses absorbable per unit, the probability that a unit is slow is
\begin{equation}
  p_{\mathrm{slow}}(d,P) =
  \begin{cases}
    1 - \displaystyle\sum_{k=0}^{K(d)} \frac{e^{-\mu}\mu^k}{k!}, & S(d) \ge 0,\\[3mm]
    1, & S(d) < 0.
  \end{cases}
  \label{eq:pslow}
\end{equation}
Delay enters twice (it shrinks $S$ and raises $r$); loss enters once
(through $\mu$).

\subsubsection{From slow units to skips}
\label{sec:model-drain}
A slow unit is not yet a skip. The decoder's lead over the content
clock evolves as
\begin{equation}
  L_{k+1} = \min\big(H,\; L_k + T_G - c_k\big),
  \label{eq:drain}
\end{equation}
and a unit is skipped once the lead has drained past one period. This
is a reflected random walk with drift $T_G - \mathbb{E}[c]$, which
yields two regimes. In the \emph{graded} regime
($\mathbb{E}[c] < T_G$), isolated slow units are absorbed; skips occur
only when a run of slow units drains the lead, so they arrive in blocks
and sessions at identical settings differ widely (the calibration
measured 6\% and 55\% at one setting). In the \emph{collapse} regime
($\mathbb{E}[c] \ge T_G$), the lead drains in about
$H/(\mathbb{E}[c] - T_G)$ units and does not recover; for pure delay
the boundary $\mathbb{E}[c] = T_G$ reduces to the threshold of
Section~\ref{sec:model-slow}. The map from $p_{\mathrm{slow}}$ to
$\tlf$ has no closed form, since it depends on $H$ and on the skip
policy. Layer~1 therefore uses Eq.~\eqref{eq:pslow} to locate the
regimes and explain the interaction, reproduces $\tlf$ by Monte Carlo
simulation of Eq.~\eqref{eq:drain}, and hands the \emph{measured}
$\tlf$ of each session to Layer~2.

\subsubsection{Non-separability and delivery-unit size}
\label{sec:model-interaction}
Equation~\eqref{eq:pslow} is a Poisson tail whose cutoff $K$ depends
on delay and whose mean $\mu$ depends on loss. Its mixed derivative
$\partial^2 p_{\mathrm{slow}}/\partial d\,\partial P$ is non-zero,
whereas a separable model, additive or multiplicative in $d$ and $P$,
requires it to vanish. The operating points of
Section~\ref{sec:design-network} make this concrete: delay alone at
\SI{9}{\milli\second} leaves $S \approx \SI{44}{\milli\second}$, below
threshold, so no misses; loss alone at the chosen rate gives
$\mu < 1$, so no misses at LAN delay; together, 50--78\% of units are
slow, where the additive prediction is zero. For $\mu \approx 0.25$ to
produce that, one loss must cost more than the whole slack;
retransmission at \SI{9}{\milli\second} does not, the window cut does,
which is why $r(d)$ is effective rather than nominal. The delivery-unit
size enters through the budget: $B_{\mathrm{eff}}$ and $S$ grow with
$G$ while $r(d)$ does not, so a larger \gof absorbs more losses per
unit. The present study fixes $G = 4$, so this prediction rests on the
earlier \gls{ra} and \gls{ldp} measurements~\cite{mosaic2026} and is
stated, not retested (Section~\ref{sec:conclusion}).

\subsection{Layer 2: Perceptual Model}
\label{sec:model-layer2}

Layer~2 maps $\tlf$ to the mean opinion score of the rated attribute.
Let $M(\tlf)$ be the expected rating on the five-point \gls{acr}
scale, with $M(0) = M_{0,c}$, the clean-reference rating of content
$c$, and $M \to 1$ as playout vanishes. The IQX hypothesis of Fiedler
et al.~\cite{fiedler2010} gives the shape: sensitivity to an
impairment is proportional to the quality headroom above the floor,
$\mathrm{d}M/\mathrm{d}\tlf = -\alpha (M - 1)$. Integrating with
$M(0) = M_{0,c}$ and normalizing the ceiling to the scale,
$\qc = (M_{0,c} - 1)/4 \in (0, 1]$,
\begin{equation}
  \crystalmodel.
  \label{eq:crystal}
\end{equation}
This is the \crystal model. The 1 and 4 are the \gls{acr} scale; $\qc$
is read off Slice~A, where $\tlf = 0$ and the equation collapses to
$1 + 4\qc$ identically,
\begin{equation}
  \qc = \frac{\mosm_{A,c} - 1}{4};
  \label{eq:ceiling}
\end{equation}
and $\alpha$ is the single fitted parameter, with the reading
$\tlf_{1/2} = \ln 2/\alpha$, the skip fraction at which perceived
smoothness has lost half its headroom. Because the clean slice is
rendered without stalls, the three $\qc$ are expected to be close to
one and statistically indistinguishable, in which case
Eq.~\eqref{eq:crystal} is a one-parameter curve shared by all content.
Objective fidelity does not enter the equation: spatial quality is
fixed within a session and has no bearing on perceived smoothness, so
a coefficient on \pcqm would only absorb variance that belongs to
content.

\emph{Freeze extension.} We also fit
$1 + 4\qc\,e^{-\alpha\tlf - \alpha_\phi \phi}$ and test
$\alpha_\phi = 0$. If freezes carry no weight beyond skips, the base
model stands; if they do, $\alpha_\phi$ is reported and Claim~1 holds
in its two-variable form either way. This test is specified before
data collection.

Compared with Eq.~\eqref{eq:udora}, which fits delay and loss directly
with per-level coefficients, \crystal moves the explanatory burden into
a measured variable and leaves one constant to be learned from people.
The two layers are falsifiable separately: Layer~1 predicts logged
quantities from settings and clean-slice telemetry and can be checked
without any rating; Layer~2 is fitted on measured $\tlf$, so an error
in Layer~1 cannot be hidden by $\alpha$.

\subsection{Predictions}
\label{sec:model-predictions}
Layer~1 makes four predictions, each checked against logs without a
rater (Table~\ref{tab:predictions}); Layer~2 is checked against
ratings. The subjective study tests Claim~1 in two forms. First, across
the impaired slices of Section~\ref{sec:design-network}, ratings fall
on one curve in $\tlf$: once $\tlf$ (and $\phi$) are accounted for,
slice identity explains no further variance. Second, the single
$\alpha$ fitted on a subset of conditions predicts held-out conditions
without refitting, at a held-out \gls{plcc} exceeding the 0.57--0.66
of Eq.~\eqref{eq:udora}. Non-separability of delay and loss is a
Layer~1 prediction: the combined slice~D is built from a delay and a
loss level that are each null on their own
(Section~\ref{sec:calibration}).

\begin{table}[t]
  \caption{Layer~1 predictions and how each is checked.}
  \label{tab:predictions}
  \centering
  \small
  \begin{tabularx}{\textwidth}{@{}>{\raggedright\arraybackslash}p{5.2cm}>{\raggedright\arraybackslash}p{3.6cm}>{\raggedright\arraybackslash}X@{}}
    \toprule
    Prediction & Form & Checked against \\
    \midrule
    Pure-delay threshold $d^\ast = B_{\mathrm{eff}} - \tau_{\mathrm{dl}}$ & closed form, no parameter & delay scan (51--55 predicted, 50--\SI{55}{\milli\second} measured) \\
    Slow-unit probability $p_{\mathrm{slow}}(d,P)$ & closed form, $r_0$ and $\kappa$ from logs & logged slow-unit fraction per session \\
    Skip fraction $\tlf$ and its between-session spread & Monte Carlo of Eq.~\eqref{eq:drain} & logged $\tlf$ per session \\
    Regime boundary $\mathbb{E}[c] = T_G$ & derived & where sessions collapse \\
    \bottomrule
  \end{tabularx}
\end{table}

The model's assumptions, in order of entry, are: loss produces no
visual artifact (reliable transport); delay is invisible while the
per-unit cycle keeps pace (buffered deadline playout); perceived
smoothness depends on the network only through $(\tlf, \phi)$
(Claim~1, tested); the IQX sensitivity law (inherited; carries the one
fitted constant); independent, Poisson-approximated loss; an affine
effective recovery cost in $d$ with measured parameters; a per-unit
budget set by the client's render frame, logged per session and held
constant across participants; and skips arising from a buffer drain
rather than from independent per-unit failures, with $\tlf$ measured
per session. Only the third is new, and only the fourth carries a
fitted constant; the last two are client properties, so the results
transfer to clients with the same frame and buffer.
